% \documentclass{acmart}
% \renewcommand\footnotetextcopyrightpermission[1]{}
% \setcopyright{none}
% \settopmatter{printacmref=false}
% \usepackage{amsmath}
% \usepackage{amsfonts}
% \usepackage[utf8]{inputenc}
% \newcommand{\x}{\mathbf{x}}
% \usepackage{float}
% \usepackage{graphicx}
% \usepackage{subcaption}

% \AtBeginDocument{%
%   \providecommand\BibTeX{{%
%     \normalfont B\kern-0.5em{\scshape i\kern-0.25em b}\kern-0.8em\TeX}}}

%  \citestyle{acmauthoryear}

% %%
% \title{Project Proposal: Machine Learning with Graphs}
% \author{Name - ID}
% \author{Name - ID}
% \author{Name - ID}
% \author{Name - ID}
% \date{}

% \begin{document}

% \maketitle
% % --------------------------------------------------
% % Project-Proposal Template  (≤ 3 pages total)
% % --------------------------------------------------

% % ---------- TL;DR ----------
% %
% % Two lines MAX that capture the essence of the project.
% %   – e.g., “We propose a new method for …” / “We extend …” / “We analyse …”.
% %
% \paragraph{TL;DR} \textit{Explain in 1–2 lines the bottom line of the project. For example: “We propose a new method for …”, “We extend …”, or “We analyse …”.}

% % ---------- 1. Background & Motivation  (≤ 1 page) ----------
% \section{Background \& Motivation}

% \textit{In <= 1 page, cover the background, motivation, and problem statement.  
% Include relevant citations. After reading this, it should be crystal-clear  
% (i) what problem you aim to solve, (ii) why it is essential, and  
% (iii) why existing approaches are insufficient (high-level).}

% %  (You may delete the bullet list once you have real prose.)
% \begin{itemize}
%   \item[•] \textbf{Problem statement} — precisely define the gap you tackle.  
%   \item[•] \textbf{Why it matters} — real-world or scientific significance.  
%   \item[•] \textbf{Prior work limits} — cite 3–6 key papers; highlight the unsolved bottleneck.
% \end{itemize}

% % ---------- 2. Proposed Approach  (≤ 1 page) ----------
% \section{Proposed Approach}

% \textit{In <= 1 page, describe how you plan to solve the problem.  
% If you propose a new method, outline it formally or verbally.  
% Be clear and concise; rough ideas are fine at this stage.  
% You may deviate slightly in the final submission if needed.}

% \begin{itemize}
%   \item[•] \textbf{Core idea} — a one-paragraph summary.  
%   \item[•] \textbf{Technical outline} — main components / algorithmic steps.  
%   \item[•] \textbf{Novelty} — what’s new relative to existing work.  

% \end{itemize}

% % ---------- 3. Experimental Plan  (≤ 1 page) ----------
% \section{Experimental Plan}

% \textit{In <= 1 page, sketch the experiments you plan to run,  
% datasets you will use, evaluation metrics, and baselines or comparison methods.  
% This is a rough plan; minor changes in the final project are allowed.}

% \begin{itemize}
%   \item[•] \textbf{Datasets} — name datasets your will use, with citations.  
%   \item[•] \textbf{Baselines / comparisons} — models or methods you’ll compare against, with citations.  
%   \item[•] \textbf{Metrics} — quantitative (and, if relevant, qualitative) measures, how you are going to measure success.  

% \end{itemize}

% \bibliographystyle{ACM-Reference-Format}
% \bibliography{references}

% \end{document}

% \documentclass{acmart}
\documentclass[manuscript,oneside]{acmart}
\makeatletter
\@twosidefalse
\makeatother
\renewcommand\footnotetextcopyrightpermission[1]{}
\setcopyright{none}
\settopmatter{printacmref=false}
\usepackage{amsmath}
\usepackage{amsfonts}
\usepackage[utf8]{inputenc}
\newcommand{\x}{\mathbf{x}}
\usepackage{float}
\usepackage{graphicx}
\usepackage{subcaption}

\AtBeginDocument{%
  \providecommand\BibTeX{{%
    \normalfont B\kern-0.5em{\scshape i\kern-0.25em b}\kern-0.8em\TeX}}}

 \citestyle{acmauthoryear}

%%
\title{Project Proposal: Machine Learning with Graphs}
\author{Yishai Lavi - 209713031}
\author{Noya Hochwald - 205420318}
\author{Omri Fahn - 206484891}
\author{Eran Kramf - 208522102}
\date{}

\begin{document}
\maketitle

\paragraph{TL;DR}
We represent a frozen ViT's attention as a graph and train a GNN to
predict failures beyond confidence and non-graph baselines.

\section{Background \& Motivation}

Vision Transformers (ViTs) represent an image as patch tokens and
repeatedly exchange information through multi-head self-attention
\citep{dosovitskiy2021image}. Reliability is
usually estimated from the final logits,
although predictions with similar confidence may arise from different
layer-wise interactions. We therefore ask:

\begin{quote}
\emph{Can a GNN operating on a ViT's multi-layer attention graph
predict classification errors, especially under unseen corruptions,
better than output-, representation-, and non-graph attention
baselines?}
\end{quote}

This is useful for selective prediction, where likely failures are
rejected or referred for review, and scientifically useful for testing
whether errors exhibit recurring routing signatures. Each image
naturally induces a directed, weighted graph over patch tokens, while
correctness is a graph-level property of the frozen classifier.

\paragraph{Related work and gap.}
ViT attention has already been interpreted as a graph, but mainly for
other objectives: relevance propagation for explanation
\citep{chefer2021transformer}, weighted-PageRank token pruning
\citep{wang2024zerotprune}, and graph regularization of class--patch
interactions for weakly supervised segmentation
\citep{yang2025more}. \citet{cultrera2023leveraging} use a
convolutional autoencoder over ViT attention heatmaps for OOD
detection, without learning on token-to-token connectivity. Closest
in task, LogitDynamics predicts ViT errors from layer-wise class-logit
trajectories, but does not use attention graphs
\citep{beigelman2026logitdynamics}.

NLP provides the closest methodological precedents. Handcrafted
topological or spectral summaries of attention graphs have been used
for generated-text and hallucination detection
\citep{kushnareva2021artificial,binkowski2025hallucination}, while
HalluZig models how attention-graph topology evolves across layers
\citep{samaga2026halluzig}. More directly, HaluGNN uses hidden states
as node features and attention values as weighted edges
\citep{kong2026halugnn}. CHARM represents raw attention scores and
activations as attributes of a token graph and applies edge-aware
message passing; its ablations show a benefit from graph connectivity
over a set model \citep{frasca2026charm}. CausalGaze further refines
attention edges using gradient sensitivity before GNN classification
\citep{kong2026causalgaze}. Thus, learned attention-graph diagnostics
are established for language models. However, our literature search
found no analogous method trained on per-image, multi-layer ViT
patch-attention graphs to predict visual classification failure. We
investigate this vision-specific gap.

\section{Proposed Approach}

Let $A^{(\ell,h)}\in\mathbb{R}^{N\times N}$ be the attention matrix at
ViT layer $\ell$ and head $h$, including the class token. Since query
token $i$ reads information from token $j$, we orient information flow as
$j\!\rightarrow\!i$. Following CHARM \citep{frasca2026charm}, the raw
per-head attention scores are used directly as edge features:
\[
\mathbf e_{ji}^{(\ell)}
=
\left[
A_{ij}^{(\ell,1)},\ldots,
A_{ij}^{(\ell,H)}
\right].
\]
To obtain a sparse graph without discarding interactions that are strong
in only one specialized head, we retain the edge $j\!\rightarrow\!i$
whenever
\[
\max_h A_{ij}^{(\ell,h)}>\tau.
\]
Thus, the threshold determines only whether the edge is present; once
retained, its edge feature contains the original attention scores from
all heads. The threshold $\tau$ is selected on validation data.

Diagonal scores are
stored as node features, together with normalized patch coordinates
and a class-token indicator. We compare this attention-focused
representation with a richer variant that also appends the
corresponding ViT token embedding $\mathbf z_i^{(\ell)}$.


An edge-aware message-passing GNN uses
$\mathbf e_{ji}^{(\ell)}$ when computing messages. We construct one
attention graph for every ViT layer,
\[
G^{(1)}(x),\ldots,G^{(L)}(x),
\]
and encode all layer graphs with a shared GNN:
\[
\mathbf g_\ell=\operatorname{GNN}\!\left(G^{(\ell)}(x)\right).
\]
The resulting ordered sequence of graph embeddings
\[
\mathbf g_1,\ldots,\mathbf g_L
\]
is then combined and followed by an MLP that predicts classification
failure. Using all layers allows the detector to capture the complete
evolution of attention routing, including early local interactions and
later global aggregation. As ablations, we
will compare this full layer-wise model with a last-layer-only model and a final-four-layers model.

For image $x$ with label $y(x)$ and frozen ViT $f$, the automatic
target is
\[
y_{\mathrm{err}}(x)
=
1\left[
\arg\max_c f_c(x)\neq y(x)
\right].
\]
Only the detector is trained, using class-weighted binary
cross-entropy. We compare:

\begin{enumerate}
    \item the attention graph with minimal node features;
    \item the attention graph augmented with token embeddings.
\end{enumerate}


\paragraph{Novelty.}
The novelty is not the attention-as-graph view or GNNs over
Transformer traces in general. It is their controlled evaluation for
per-image ViT patch graphs, with comparisons against output signals,
hidden representations, and non-graph models given the same attention
data.

\section{Experimental Plan}

\paragraph{Datasets and setup.}
The primary benchmark will use CIFAR-100 and CIFAR-100-C, whose
corruption families and severity levels follow
\citet{hendrycks2019benchmarking}. We will begin with a publicly available ViT checkpoint fine-tuned on
CIFAR-100 and verify its accuracy, preprocessing pipeline, label
mapping, confidence distribution, and access to attention matrices from
all Transformer layers. If no suitable checkpoint is available, we
will fine-tune an ImageNet-pretrained ViT on CIFAR-100. The selected
ViT will then be frozen throughout all detector experiments.

If no publicly available checkpoint satisfies these requirements, we
will fine-tune an ImageNet-pretrained ViT on the clean CIFAR-100 training
set and select the checkpoint according to validation accuracy. In
either case, the selected ViT will then be frozen and used as the fixed
classifier throughout all detector experiments.

We will evaluate the selected classifier on CIFAR-100-C and report
accuracy separately for each corruption type and severity level. After
validating the classifier, we will precompute its attention graphs. The
detector will be trained on clean images and a selected subset of noise,
blur, weather, and digital corruptions, while complete corruption
families will be reserved for testing. All corrupted versions of the
same source image will remain in the same split to prevent leakage.
Optional far-OOD datasets are SVHN \citep{netzer2011svhn} and DTD
\citep{cimpoi2014describing}; localized patches and occlusions will be
generated automatically.


\paragraph{Baselines and ablations.}
We first report the frozen ViT's clean and corruption-specific
classification accuracy. For failure prediction, we compare against
maximum softmax probability, logit margin, and an MLP on the final
class-token embedding.

To isolate the contribution of graph structure, we also compare against
an MLP receiving the same flattened attention tensors as the GNN and a
sequence model receiving class-token embeddings from all ViT layers.

Graph ablations include last-layer-only input, the final four layers,
all layers, and threshold sensitivity.

\paragraph{Metrics and success criteria.}
We evaluate error prediction using AUROC and AUPRC, treating an
incorrect ViT prediction as the positive class. Since the detector may
be used to reject unreliable predictions, we also report
risk--coverage curves, area under the risk--coverage curve, and
selective accuracy at fixed coverage levels.

Metrics are reported separately for clean images, corruption families
seen during detector training, and unseen corruption families. The main
hypothesis is supported if the GNN outperforms non-graph models given
the same attention values, particularly on unseen corruptions.

\paragraph{Compute resources.}
The project requires GPU access provided by the course or university. The
main compute cost is fine-tuning the ViT and precomputing attention graphs,
token features, and logits for clean and corrupted images, as well as training the GNN over them.

\bibliographystyle{ACM-Reference-Format}
\bibliography{references}
\end{document}
